% ==============================================================
%  Chapter 4 -- Experimental Study
% ==============================================================

\chapter{Experimental Study}
\label{ch:experiments}

This chapter puts the pipeline of Chapter~\ref{ch:implementation} to
the test on two real influence diagrams of clinical origin. The
\emph{bypass} diagram, already introduced in
Section~\ref{sec:bypass-example} as the running example of the
implementation, plays the role of a small sanity-check benchmark
where the parameter recovery can be tracked node by node. The
\emph{NHLv1} diagram, presented in the next section, is one order of
magnitude larger and constitutes the main benchmark of the empirical
study.

% ==============================================================
\section{Benchmark networks}
\label{sec:benchmarks}
% ==============================================================

\subsection{NHLv1: management of non-Hodgkin lymphoma}
\label{sec:nhlv1}

The \emph{NHLv1} influence diagram models the clinical management of
gastric non-Hodgkin lymphoma. Its qualitative structure is shown in
Figure~\ref{fig:nhlv1-id}. It comprises sixteen chance nodes
covering disease, treatment-side-effect and outcome variables, three
decision nodes (\texttt{Surgery}, \texttt{Helicobacter treatment} and
\texttt{CT-RT schedule}) and a single utility node. Because of its
size and connectivity, the parameter vector
$\boldsymbol{\Theta}$ associated with the diagram is of the order of
several thousand entries, which makes the elicitation bottleneck
discussed in Chapter~\ref{ch:introduction} fully tangible and the
EDA-based recovery the only practical route to populating the
network.

\begin{landscape}
\begin{figure}[p]
\centering
\begin{tikzpicture}[
    x=1.35cm, y=1.25cm, % Amplía el espacio entre nodos para evitar solapamientos
    every node/.style={font=\sffamily\scriptsize, align=center, inner sep=2pt},
    chance/.style={ellipse, draw=blue!50, fill=cyan!15,
                   minimum width=2.2cm, minimum height=0.9cm},
    decision/.style={rectangle, draw=black, fill=orange!30,
                     minimum width=2.2cm, minimum height=1.0cm},
    utility/.style={regular polygon, regular polygon sides=6,
                    shape border rotate=90,
                    draw=black, fill=purple!30,
                    minimum width=2.4cm, minimum height=1.0cm, inner sep=2pt},
    arr/.style={-{Stealth[scale=1.0]}, draw=blue!60!black},
    info/.style={-{Stealth[scale=1.0]}, draw=blue!60!black, dashed}
]
    % ==========================================
    %  Posición de Nodos (Ajustados al Layout)
    % ==========================================

    %  --- Fuentes / observaciones ---
    \node[chance]   (AGE)    at ( 0, 10)   {AGE};
    \node[chance]   (GHS)    at ( 0,  6)   {GENERAL HEALTH\\STATUS};
    \node[chance]   (HPYL)   at ( 0, -1)   {HELICOBACTER\\PYLORI};

    %  --- Decisiones ---
    \node[decision] (SURG)   at ( 3.5,  5)   {SURGERY};
    \node[decision] (HTRT)   at ( 2,  1)   {HELICOBACTER\\TREATMENT};
    \node[decision] (CTRT)   at ( 4, 10)   {CT RT\\SCHEDULE};

    %  --- Erradicación ---
    \node[chance]   (ERAD)   at ( 3, -1)   {ERADICATION};

    %  --- Cadena Quirúrgica y Clínica ---
    \node[chance]   (PSRV)   at ( 6,  7.5) {POST SURGICAL\\SURVIVAL};
    \node[chance]   (CLIN)   at ( 6,  4)   {CLINICAL\\PRESENTATION};
    \node[utility]  (UTIL)   at ( 6,  1.5) {UTILITY};
    \node[chance]   (HIST)   at ( 6, -1)   {HISTOLOGICAL\\CLASSIFICATION};

    %  --- Cadena CT-RT ---
    \node[chance]   (HEM)    at ( 9, 10)   {HEMORRHAGE};
    \node[chance]   (BM)     at ( 8,  9)   {BM DEPRESSION};
    \node[chance]   (BULKY)  at ( 9.5,  2.5)   {BULKY DISEASE};

    \node[chance]   (IMM)    at (8,  6)   {IMMEDIATE\\SURVIVAL};
    \node[chance]   (POCT)   at (11,  8)   {POST CT RT\\SURVIVAL};
    \node[chance]   (CSTG)   at (11, 0)   {CLINICAL STAGE};

    \node[chance]   (PERF)   at (13,  6)   {PERFORATION};
    \node[chance]   (THER)   at (13, 10)   {THERAPY\\ADJUSTMENT};

    %  --- Endpoints ---
    \node[chance]   (EARLY)  at (15,  4)   {EARLY RESULT};
    \node[chance]   (FIVE)   at (15,  0)   {FIVE YEAR\\RESULT};


    % ==========================================
    %  Arcos (Totalmente Rectos)
    % ==========================================
    
    %  AGE -> ...
    \draw[arr] (AGE) -- (GHS);
    \draw[arr] (AGE) -- (BM);
    \draw[arr] (AGE) -- (HIST);
    \draw[arr] (AGE) -- (CSTG);
    \draw[arr] (AGE) -- (BULKY);
    \draw[arr] (AGE) -- (FIVE);

    %  GHS / CLIN / SURG -> PSRV
    \draw[arr] (GHS)  -- (PSRV);
    \draw[arr] (CLIN) -- (PSRV);
    \draw[arr] (SURG) -- (PSRV);

    %  CTRT -> ...
    \draw[arr] (CTRT) -- (HEM);
    \draw[arr] (CTRT) -- (BM);
    \draw[arr] (CTRT) -- (PERF);
    \draw[arr] (CTRT) -- (EARLY);

    %  Complicaciones -> THER / POCT
    \draw[arr] (BM)   -- (THER);
    \draw[arr] (HEM)  -- (THER);
    \draw[arr] (PERF) -- (THER);
    \draw[arr] (HEM)  -- (POCT);
    \draw[arr] (PERF) -- (POCT);

    %  Survival chain
    \draw[arr] (POCT) -- (IMM);
    \draw[arr] (PSRV) -- (IMM);

    %  Surgery / decisions -> outcomes
    \draw[arr] (SURG) -- (EARLY);
    \draw[arr] (THER) -- (EARLY);

    %  HTRT chain
    \draw[arr] (HPYL) -- (ERAD);
    \draw[arr] (HTRT) -- (ERAD);
    \draw[arr] (ERAD) -- (EARLY);

    %  Classification / staging -> outcomes
    \draw[arr] (HIST) -- (BULKY);
    \draw[arr] (CSTG) -- (BULKY);
    \draw[arr] (BULKY) -- (EARLY);
    \draw[arr] (HIST) -- (EARLY);
    \draw[arr] (CSTG) -- (EARLY);

    %  EARLY -> FIVE and parents of FIVE
    \draw[arr] (EARLY) -- (FIVE);
    \draw[arr] (CSTG)  -- (FIVE);
    \draw[arr] (BULKY) -- (FIVE);
    \draw[arr] (HIST)  -- (FIVE);
    \draw[arr] (IMM)   -- (FIVE);

    %  Utility parents
    \draw[arr] (HTRT) -- (UTIL);
    \draw[arr] (SURG) -- (UTIL);
    \draw[arr] (CTRT) -- (UTIL);
    \draw[arr] (GHS)  -- (UTIL);
    \draw[arr] (FIVE) -- (UTIL);


    % ==========================================
    %  Arcos informacionales (Líneas Discontinuas)
    % ==========================================
    
    %  Hacia HTRT
    \draw[info] (HPYL)  -- (HTRT);
    \draw[info] (CSTG)  -- (HTRT);
    \draw[info] (HIST)  -- (HTRT);
    \draw[info] (BULKY) -- (HTRT);
    
    %  Hacia SURG
    \draw[info] (GHS)  -- (SURG);
    \draw[info] (CLIN) -- (SURG);
    \draw[info] (HTRT) -- (SURG);
    
    %  Hacia CTRT
    \draw[info] (SURG) -- (CTRT);

\end{tikzpicture}
\caption{Influence diagram of the \emph{NHLv1} model for the management of gastric non-Hodgkin lymphoma.}
\label{fig:nhlv1-id}
\end{figure}
\end{landscape}

\subsection{Bypass: management of coronary heart disease}
\label{sec:bypass-bench}

The \emph{bypass} diagram was already presented in
Section~\ref{sec:bypass-example} as the illustrative model used to
explain the implementation. It is reused here as a small, controlled
benchmark on which the parameter vector is small enough for a finer
ablation study of the encoding, the decoding, the loss and the EDA
variant.

% ==============================================================
\section{Experimental protocol}
\label{sec:protocol}
% ==============================================================

% TODO: describe the hyperparameter sweep, the metrics
% (rule accuracy, MSE w.r.t. true CPTs, MSE w.r.t. true utility,
% utility deviation, normalised entropy, training/validation split,
% number of repetitions, random seeds, hardware).

\bigskip
\textbf{[To be completed.]}

% ==============================================================
\section{Results}
\label{sec:results}
% ==============================================================

% TODO: present the results grouped by axis of the sweep:
%  - EDA variant (UMDA vs EMNA vs EGNA vs KEDA)
%  - loss function (binary / regret / margin / softmax / entropy)
%  - temperatures (chance / utility)
%  - selection ratio alpha and population size
%  - symmetric vs non-symmetric sampling
%  - min_max_ut on/off
%  - mode (both / chance_only / utility_only)
%  - effect of rule-base size

\bigskip
\textbf{[To be completed.]}

% ==============================================================
\section{Discussion}
\label{sec:discussion}
% ==============================================================

% TODO: summarise the empirical findings, link them back to the
% theoretical analysis of Chapter 2 (non-convexity, degeneracy),
% and highlight the limits.

\bigskip
\textbf{[To be completed.]}
